% Slides — Capítulo 20: Previsão Numérica do Tempo
\documentclass[aspectratio=169,11pt]{beamer}
\usepackage{../comum/temameteo}
\graphicspath{{../figuras/cap20/}}

\title[Cap. 20 --- Previsão Numérica]{Capítulo 20 --- Previsão Numérica
do Tempo}
\subtitle{Meteorologia Prática}
\author{}
\date{}

\begin{document}

\begin{frame}
  \titlepage
  \vfill
  \atribuicao
\end{frame}

\begin{frame}{Roteiro}
  \tableofcontents
\end{frame}

% ---------------------------------------------------------------
\section{As equações e a grade}

\begin{frame}{O globo em caixinhas}
  \begin{columns}
    \column{0.3\textwidth}
    \begin{itemize}
      \item As 7 equações do curso (Caps.~1, 3, 4, 10) em $10^9$
            células.
      \item Grades esféricas, malhas aninhadas, coordenada $\sigma$
            seguindo o terreno.
      \item Arakawa C: cada variável no seu canto da célula.
    \end{itemize}
    \column{0.35\textwidth}
    \figlivroslide{fig_20_5a.jpg}{0.42\textheight}{Fig.~20.5 --- a
      esfera cubada de um modelo global}
    \column{0.35\textwidth}
    \figlivroslide{fig_20_2a.jpg}{0.42\textheight}{Fig.~20.2 ---
      coordenada $\sigma$ abraçando a montanha}
  \end{columns}
\end{frame}

\begin{frame}{Advecção discreta: os pecados de cada esquema}
  \begin{columns}
    \column{0.4\textwidth}
    \eqdestaque{CFL: $C = \dfrac{U\Delta t}{\Delta x} \le 1$}
    \vspace{0.4em}
    \begin{itemize}
      \item Upwind difunde ($K_{num} = U\Delta x(1-C)/2$, T2);
            leapfrog dispersa; $C>1$ explode (T1; Caso~1).
      \item A onda $2\Delta x$ viaja errada: resolução efetiva
            $\sim 7\Delta x$ (T5).
    \end{itemize}
    \column{0.6\textwidth}
    \figlivroslide{fig_20_11c.jpg}{0.5\textheight}{Fig.~20.11 ---
      violando o CFL: a solução explode}
  \end{columns}
\end{frame}

\begin{frame}{Parametrizações: a física sub-grade}
  \begin{columns}
    \column{0.4\textwidth}
    \begin{itemize}
      \item O que a célula não vê vira estatística: radiação
            (Cap.~2), microfísica (Cap.~7), convecção (Cap.~14),
            turbulência (Cap.~18), ondas de montanha (Cap.~17).
      \item Metade do custo computacional — e a maior fonte de erro
            sistemático.
      \item Entender os Caps.~2--18 $=$ saber onde o modelo erra.
    \end{itemize}
    \column{0.6\textwidth}
    \figlivroslide{fig_20_3c.jpg}{0.48\textheight}{Fig.~20.3 --- os
      processos que uma célula precisa parametrizar}
  \end{columns}
\end{frame}

% ---------------------------------------------------------------
\section{Caos e previsibilidade}

\begin{frame}{Lorenz 1963: a borboleta}
  \begin{columns}
    \column{0.4\textwidth}
    \eqdestaque{$t_{max} = \dfrac{1}{\lambda}
      \ln\dfrac{E}{\epsilon_0}$}
    \vspace{0.4em}
    \begin{itemize}
      \item Gêmeos indistinguíveis divergem exponencialmente
            (Caso~2; N2).
      \item Logarítmico em $\epsilon_0$: melhorar observações 10×
            compra só $\sim$2 dias (T3).
      \item Teto de $\sim$2 semanas: estrutural, não computacional.
    \end{itemize}
    \column{0.6\textwidth}
    \figlivroslide{fig_20_18a.jpg}{0.52\textheight}{Fig.~20.18 --- o
      atrator do sistema de Lorenz}
  \end{columns}
\end{frame}

% ---------------------------------------------------------------
\section{Assimilação e ensembles}

\begin{frame}{A média ponderada mais valiosa do mundo}
  \eqdestaque{$x_a = x_b + K(y_o - x_b)$,\qquad
    $K^* = \dfrac{\sigma_b^2}{\sigma_b^2+\sigma_o^2}$,\qquad
    $\dfrac{1}{\sigma_a^2} = \dfrac{1}{\sigma_b^2} +
    \dfrac{1}{\sigma_o^2}$}
  \vspace{0.5em}
  \begin{itemize}
    \item As precisões SOMAM: a análise supera modelo E observação
          (T4; Caso~3; N3).
    \item Ciclada a cada 6\,h sobre $10^7$ observações (Caps.~8--9):
          satélites são a espinha dorsal.
    \item Generalizações: 4D-Var, EnKF — o coração é esta fórmula (a
          mesma do filtro de Kalman e do laboratório de física).
  \end{itemize}
\end{frame}

\begin{frame}{Ensemble: prever a própria incerteza}
  \begin{columns}
    \column{0.4\textwidth}
    \begin{itemize}
      \item 30--50 rodadas perturbadas: a dispersão é a barra de
            erro (Caso~4).
      \item Fração de membros $=$ probabilidade: ``70\% de chuva'' é
            saída literal.
      \item O spread prevê o próprio erro (N4): o ensemble sabe
            quando vai errar.
    \end{itemize}
    \column{0.6\textwidth}
    \figlivroref{Fig.~20.19}{o espaguete do ensemble abrindo com o prazo}{ECMWF}
  \end{columns}
\end{frame}

\begin{frame}{Verificação: a revolução silenciosa}
  \begin{columns}
    \column{0.4\textwidth}
    \begin{itemize}
      \item Previsão de 5 dias hoje $=$ 3 dias dos anos 1990:
            $\sim$1 dia de ganho por década.
      \item Ingredientes: satélites $+$ assimilação $+$ física $+$
            computação.
      \item Persistência e climatologia: os adversários que toda
            previsão precisa vencer.
    \end{itemize}
    \column{0.6\textwidth}
    \figlivroslide{fig_20_22.jpg}{0.5\textheight}{Fig.~20.22 ---
      correlação de anomalia vs.\ prazo: previsão × persistência}
  \end{columns}
\end{frame}

% ---------------------------------------------------------------
\section{Síntese e laboratório}

\begin{frame}{O mapa do capítulo}
  \eqdestaque{$C \le 1$ \quad $K_{num} = \tfrac{U\Delta x(1-C)}{2}$
    \quad $t_{max} = \tfrac{1}{\lambda}\ln\tfrac{E}{\epsilon_0}$
    \quad $x_a = x_b + K(y_o-x_b)$}
  \vspace{0.6em}
  Discretizar, parametrizar, inicializar, quantificar: o curso inteiro
  dentro de um computador — com prazo de validade dado pelo caos e
  incerteza entregue pelo ensemble.
\end{frame}

\begin{frame}{Laboratório numérico --- notebook do Cap.~20}
  \texttt{notebooks/cap20\_previsao.ipynb}
  \begin{enumerate}
    \item Upwind / Lax--Wendroff / leapfrog e o CFL.
    \item Lorenz 63: borboleta e gêmeos divergindo.
    \item Ciclo de assimilação no Lorenz.
    \item Ensemble: espaguete $\to$ probabilidade.
  \end{enumerate}
  \vspace{0.4em}
  \textbf{Exercícios}: T1--T5 e N1--N4 nas notas; células reservadas no
  notebook. Rodar o Caso~2 ao vivo mudando o 6º decimal é o momento
  ``uau'' do semestre.
  \vfill
  \atribuicao
\end{frame}

\end{document}
